\section{Interacción multiturno y diseño de constitución}

El segundo paso de Compound AI es la recolección de datos mediante diálogo multiturno. Cuando el usuario solo da una intención vaga, el sistema necesita preguntar por los detalles, mantener la credibilidad y producir instrucciones estructuradas dentro de un número limitado de turnos.

\subsection{APAC Constitution}

Abstrajeron esta necesidad en la APAC Constitution: las cuatro dimensiones Accuracy, Proactivity, Adaptivity y Credibility, cada una cuantificada mediante conductas concretas, por ejemplo si las preguntas se concentran en la información realmente clave o si se evitan preguntas de seguimiento irrelevantes repetidas.
\begin{figure}[H]
\centering
\includegraphics[width=\textwidth]{frame-apac-constitution.jpg}
\caption{La APAC Constitution responde al usuario en varias dimensiones, logrando preguntas proactivas y alta confianza. \protect\footnotemark}
\end{figure}
\footnotetext{Intervalo de tiempo en el video: 00:16:08--00:16:22.}
\begin{importantbox}{Lecciones de la constitución APAC}
1) Accuracy: los detalles deben corresponder a los campos clave que el usuario necesita; 2) Proactivity: plantear proactivamente preguntas de valor; 3) Adaptivity: ajustar el estilo según distintos traveler persona; 4) Credibility: controlar la hallucination y evitar contradicciones evidentes. APAC hace posible optimizar estos objetivos de manera simultánea.
\end{importantbox}
\textit{``Agent Constitution makes it possible to improve the agent performance along different axes simultaneously with very simple techniques.''}

Un ejemplo real de diálogo es: cuando los travelers dicen ``I want to go to Hawaii'', el agent primero confirma el presupuesto, las preferencias (money vs experience) y si desea múltiples ciudades, después pregunta por la ventana de tiempo y los medios de transporte aceptables, y en cada turno actualiza la información recabada en el Json schema. Este pipeline de preguntas graduales y verificación inmediata está impulsado por los cuatro objetivos de APAC, lo que garantiza que las preguntas tengan valor sin desviarse de la intención inicial del usuario.

\subsection{Preguntas proactivas y evaluación}

Simularon 50 traveler persona, cada uno con preferencias distintas y una hidden spreadsheet; el agent completa la entrada Json mediante diálogo multiturno y después se aplica DPO fine-tune para optimizar la alineación con el ground truth. El experimento también usa al propio agent como judge, sin depender ya de una calificación humana durante todo el proceso, lo que aumenta la eficiencia de los datos.
\begin{knowledgebox}{Métricas de preguntas y respuestas multiturno}
Las evaluaciones clave incluyen: 1) overall recall (si se recabó toda la información); 2) critical recall (la corrección de los campos más importantes); 3) agentic score (la proactividad de la respuesta y la coherencia a lo largo de los turnos); 4) DPO reward (la alignment con el judge).
\end{knowledgebox}
\begin{warningbox}{Evitar la trampa de un persona simplista}
Las preguntas proactivas deben subordinarse al objetivo: preguntar en exceso desperdicia turnos, y una plantilla demasiado rígida tiende a ignorar las particularidades de cada persona. El modelo debe encontrar un equilibrio entre accuracy y efficiency.
\end{warningbox}
\begin{knowledgebox}{Heurísticas para las preguntas proactivas}
1) Rastrear la información faltante: en cada turno, ordenar las preguntas candidatas según la importancia del campo; 2) Medir el cost-benefit: si una pregunta aporta mucho valor pero tiene un costo alto, puede dividirse en dos preguntas más pequeñas; 3) Usar la persona para enfocar las diferencias: con un perfil sensible al presupuesto, preguntar rápido por el costo; con uno orientado a la experiencia, enfatizar las actividades; 4) Limitar los turnos: si se superan los 4 turnos, considerar resumir la información ya conocida.
\end{knowledgebox}

\subsection{Resumen de la sección}

La APAC Constitution dota a las preguntas y respuestas multiturno de controlabilidad y medibilidad; agent judge + DPO forman un ciclo cerrado, lo que hace que el fine-tuning autosupervisado sea un esquema de entrenamiento viable.

